% PaperLab venue template: MethodsX (methodological protocol article)
% Minimal compilable skeleton — article class, common packages only.
\documentclass[11pt]{article}
\usepackage[a4paper,margin=1in]{geometry}
\usepackage{graphicx}
\usepackage{booktabs}
\usepackage{natbib}
\usepackage{hyperref}

\title{TITLE: A Reproducible Protocol for METHOD on TASK}
\author{AUTHOR NAME\thanks{Affiliation; email@example.com}}
\date{\today}

\begin{document}
\maketitle

\begin{abstract}
\noindent PLACEHOLDER-ABSTRACT: State the methodological problem, what this protocol delivers,
the inputs it consumes, the outputs it produces, and how each step is verified.
MethodsX articles describe a method that is new or substantially improved,
with enough detail for another lab to reproduce it without specialist knowledge.
\\[2mm]
\textbf{Keywords:} keyword one; keyword two; keyword three
\end{abstract}

\section{Background}
PLACEHOLDER-BACKGROUND: why the method is needed and what current approaches lack.

\section{Method}
PLACEHOLDER-METHOD: the protocol itself, step by step. Reference figures as
Figure~\ref{fig:placeholder} and tables as Table~\ref{tab:placeholder}.

\begin{figure}[t]
\centering
% \includegraphics[width=0.9\linewidth]{figures/figure1.pdf}
\caption{PLACEHOLDER-CAPTION: what the reader should see.}
\label{fig:placeholder}
\end{figure}

\begin{table}[t]
\centering
\caption{PLACEHOLDER-TABLE: tool/parameter matrix.}
\label{tab:placeholder}
\begin{tabular}{lll}
\toprule
Step & Tool & Parameter \\
\midrule
1 & PLACEHOLDER & PLACEHOLDER \\
2 & PLACEHOLDER & PLACEHOLDER \\
\bottomrule
\end{tabular}
\end{table}

\section{Validation}
PLACEHOLDER-VALIDATION: how correctness is checked (metrics, ground truth, seeds)
and the exact numbers the validation produced.

\section{Conclusions}
PLACEHOLDER-CONCLUSIONS: what the protocol adds; limitations; where it should not be used.

\section*{Declarations}
PLACEHOLDER-DECLARATIONS: funding, conflicts of interest, data availability.

\bibliographystyle{plainnat}
\bibliography{refs}

\end{document}
